先考虑较粗的处理：四大家鱼合为单一消费者功能群，水草、浮游植物、浮游动物和底栖饵料只保留一个资源总量。这样的模型能够匹配2025年资源量恢复至1.8倍，也能验证"鱼类资源量增加"，但它无法判断哪一层级资源率先承压，题目中的"水下荒漠"等底层生态退化现象也没有相应的解释位置。问题一还要追查"鱼类增长的驱动机制"和"基础资源的承载压力变化"，因此在正式模型中保留四类基础资源与四大家鱼的基本食性对应关系。
